\begin{table}[htb!]
\centering
\footnotesize
\begin{tabular}{@{}lrrrr@{}}
\toprule
\textbf{Pair} & \textbf{Market} & \textbf{Model} & \textbf{Excess} & \textbf{Positive Years} \\
\midrule
AUD/USD & $-18.2\%$ & $+87.2\%$ & $+105.4$ pp & 8 of 11 \\
EUR/USD & $-2.9\%$ & $+70.6\%$ & $+73.4$ pp & 5 of 11 \\
GBP/USD & $-13.5\%$ & $+129.1\%$ & $+142.5$ pp & 7 of 11 \\
NZD/USD & $-26.1\%$ & $+60.1\%$ & $+86.1$ pp & 7 of 11 \\
USD/CAD & $+18.1\%$ & $+221.0\%$ & $+202.9$ pp & 8 of 11 \\
USD/CHF & $-20.3\%$ & $+54.8\%$ & $+75.2$ pp & \textbf{10 of 11} \\
USD/JPY & $+30.7\%$ & $+16.2\%$ & $\mathbf{-14.4}$ pp & 5 of 11 \\
\bottomrule
\end{tabular}
\caption{Each agent against buying and holding its own pair over the same window and under the same costs. Five of the seven pairs lost value over the eleven years, so on those the return was earned against a falling market. USD/JPY is the only pair whose benchmark rose substantially, and the only one the agent failed to beat.}
\label{Tables:Benchmark}
\end{table}